% Supplementary source and checking details; not an additional theorem.
\section{Lean formalization and verification}\label{sec:lean}

The project \texttt{lean/} uses Lean 4 \cite{lean} and mathlib
\cite{mathlib}. The complete chain from the literal greedy rule
\eqref{eq:rule}, through the smoothing \eqref{eq:smooth}, to
Theorem~\ref{thm:main} is machine-checked. Haver's estimates in
Proposition~\ref{prop:inputs} are \emph{imported proved theorems from Haver's
Lean sources}, recompiled in this development. They are neither
independently re-proved mathematical contributions of this paper nor
hypotheses or added axioms. The archive provides, among others,
\path{EKG/Canonical.lean}, \path{EKG/EarlyEvens.lean},
\path{EKG/Span.lean}, \path{EKG/MissingComposites.lean} and
\path{EKG/SmoothedPositions.lean}. Our new proofs of the two little-$o$
estimates, the endpoint hypotheses of the counting identity, and the final
transfer close the chain.
The elementary constants used to weaken Haver's bounds are also proved.

The type transliterated beside Theorem~\ref{thm:main} is independently fixed by
\texttt{EndpointStatement} in \path{Statement.lean}.
\path{Check.lean} requires \texttt{ekg\_sharp\_all\_index} to have that type.
For the rule, \texttt{actualA\_properties}, \texttt{actualA\_gcd\_rule}
and \texttt{actualA\_gcd\_minimal} prove positivity, non-repetition,
gcd eligibility, and leastness, respectively. Thus the internal
\texttt{GreedySeq} interface is instantiated by the recursively defined
sequence, rather than left as an assumption.
The proof sources accompany this manuscript in the same deposit.

The formal $a(n)$ is the zero-based canonical sequence evaluated at
$n-1$; its value at $n=0$ is a total extension unused in the theorem.
The canonical construction proves existence, minimality, eligibility,
and distinctness for the literal rule \eqref{eq:rule}, including the
prime factor $2$. The new \texttt{smoothValue} definition excludes
$2$ in both branches. The bridge to Haver's smoothing is proved
for every $n\geq5$, using that the exceptional raw values $2,6$
occur only at indices $2,4$.

The files under \texttt{EKG/} correspond to the proof as follows:
\begin{center}\small
\begin{tabular}{ll}
\texttt{PartA.lean}, \texttt{PartB.lean} & upper and lower excursions\\
\texttt{CountingExact.lean} & exact counting identity\\
\texttt{EndpointRemainders.lean} & negligible remainders\\
\texttt{EndpointClock.lean} & prime counts and frontier clock\\
\texttt{EndpointStructure.lean} & ordinary terms and prime blocks\\
\texttt{EndpointTransfer.lean} & transfer to all indices
\end{tabular}
\end{center}
The weak prime number theorem is obtained from
\texttt{pi\_alt'} in the pinned \emph{PrimeNumberTheoremAnd}
project \cite{pntcode}; its real-floor inclusive prime count is
adapted in \path{SolutionPNT.lean}.

The Lean toolchain and dependency revisions are pinned in
\texttt{lean/lean-toolchain} and \texttt{lean/lake-manifest.json}.
Source provenance is recorded in
\path{THIRD_PARTY.md}.

One inherited finite \texttt{native\_decide} proof is replaced by
ordinary kernel proofs; two compatibility edits change tactic lines
only, without changing definitions or declaration types.

The exact output of \texttt{\#print axioms ekg\_sharp\_all\_index} is
\begingroup\small
\begin{verbatim}
'ekg_sharp_all_index' depends on axioms: [propext, Classical.choice, Quot.sound]
\end{verbatim}
\endgroup
These are the usual Lean logical axioms, with no mathematical
hypothesis added for an estimate, the counting identity or the prime
number theorem.
The imported PNT file contains unrelated admitted declarations
\texttt{prelim\_decay\_2}, \texttt{prelim\_decay\_3} and \texttt{decay\_alt};
none occurs in the proof of Theorem~\ref{thm:main}. This does not certify the whole
upstream project. The three-branch corollary is the elementary algebraic
consequence proved in the text; its separate statement is not included
as an additional Lean theorem.

The following consequences in Section~\ref{sec:consequences} are also
machine-checked. \path{CorollaryStatements.lean} fixes their statements
independently of the proofs, for the one-based sequence \texttt{a}:
\begin{itemize}
\item \texttt{EventualSignStatement}: the first part of
  Corollary~\ref{cor:sign};
\item \texttt{BelowDiagonalCountStatement}: its count
  $\#\{k\leq N:a(k)<k\}\sim N/(2\log N)$;
\item \texttt{OrdowskiDensityStatement}: its last statement,
  $\#\{k\leq N:a(k)>k\}/N\to1$. The sharper form
  $N-(1/2+o(1))N/\log N$ follows from the previous count and the
  finiteness of $\{k:a(k)=k\}$ as in the proof, but is not stated
  as a separate Lean theorem;
\item \texttt{JasinskiFinitenessStatement}: Corollary~\ref{cor:lines},
  with the lines written $2a(k)=k$, $a(k)=k$ and $2a(k)=3k$ as in the OEIS;
\item \texttt{ExceptionalFinitenessStatement}: the finiteness of the
  indices $n$ with $a(n)$ not prime and $a(n)\leq n-1$, and of A064420;
\item \texttt{A155549FinitenessStatement}: the finiteness of A155549,
  the image of that index set under $a$;
\item \texttt{ShortCyclesStatement}: Corollary~\ref{cor:cycles}, a cycle
  being the orbit of a positive $m$ with $a^{[l]}(m)=m$ for some
  $1\leq l\leq L$.
\end{itemize}
The proofs are in \path{SolutionSign.lean},
\path{SolutionDensity.lean} and \path{SolutionCycles.lean}; they use
Theorem~\ref{thm:main}, the permutation property and the prime number
theorem above. \path{Check.lean} checks each against its statement and
guards its axioms, which are again exactly
\texttt{propext}, \texttt{Classical.choice} and \texttt{Quot.sound}.

No human expert has reviewed it.


Haver proves this counting identity in Corollary 12(b).
Haver's Lean archive already proves \path{clock_balance},
\path{clock_identity}, \path{clock_identity_with_ticks}, and
\path{clock_identity_from_interval_counts}; the last takes the
endpoint errors as hypotheses. We discharge those hypotheses by proving
that the missing-prime count is $P_1+\theta_1$ and the tick count is
$P_2-\theta_2$, and by evaluating the interval endpoints through $I$
and $\theta$ as above.
